\section{Critérios de sustentabilidade e priorização}
\label{sec:criterios}

No núcleo final de 104 registros, a leitura estruturada de título, resumo e palavras-chave confirmou a presença de seis dimensões de sustentabilidade em proporção relevante: técnica-operacional em 102 registros (98,1\%), institucional em 93 (89,4\%), ambiental em 84 (80,8\%), ciclo de vida em 60 (57,7\%), econômica em 59 (56,7\%) e social em 57 (54,8\%) (Tabela~\ref{tab:dimensoes}, Figura~\ref{fig:dimensoes}). A predominância da dimensão técnica-operacional é consistente com a própria natureza do objeto -- manutenção e gestão de ativos construídos --, ao passo que a presença elevada da dimensão institucional sugere que grande parte da literatura associa a manutenção predial a estruturas de governança, conformidade normativa e capacidade de execução, não apenas a procedimentos técnicos isolados.

\begin{table}[H]
\centering
\caption{Dimensões de sustentabilidade identificadas no núcleo final}
\label{tab:dimensoes}
\small
\begin{tabular}{l r r}
\toprule
Dimensão & Registros & \% dos 104 \\
\midrule
Técnica-operacional & 102 & 98,1 \\
Institucional & 93 & 89,4 \\
Ambiental & 84 & 80,8 \\
Ciclo de vida & 60 & 57,7 \\
Econômica & 59 & 56,7 \\
Social & 57 & 54,8 \\
\bottomrule
\end{tabular}

\vspace{2pt}
\begin{minipage}{0.75\linewidth}
\footnotesize Fonte: elaborado pelo autor a partir de \texttt{fontes/tabela27\_dimensoes\_sustentabilidade\_nucleo\_final\_104.csv} (seis dimensões canônicas; ver nota de qualidade de dado na Seção~\ref{sec:limitacoes}).
\end{minipage}
\end{table}

\begin{figure}[H]
\centering
\includegraphics[width=0.75\linewidth]{figuras/figura12_dimensoes_sustentabilidade_nucleo_final_104.png}
\caption{Dimensões de sustentabilidade identificadas no núcleo final}
\label{fig:dimensoes}
\end{figure}

Entre os critérios de priorização identificados, o desempenho operacional aparece em 93 registros (89,4\%), seguido da disponibilidade de dados e informação em 76 (73,1\%), do custo em 59 (56,7\%), da vida útil do ativo em 40 (38,5\%) e da energia em 36 (34,6\%) (Tabela~\ref{tab:criteriospriorizacao}). Critérios de condição física (27 registros) e de risco (31 registros) aparecem como recortes técnico-operacionais explícitos, ao passo que critérios sociais -- conforto, segurança e satisfação do usuário -- e critérios ambientais específicos -- emissões de carbono, resíduos e consumo de água -- aparecem com frequência mais baixa, mas de forma textualmente explícita nos resumos analisados. \textcite{kim_estimativacustoseducacional_2018}, ao desenvolverem um modelo de estimativa de custo de manutenção e reparo para edificações educacionais por regressão, ilustram exatamente esse tipo de critério econômico-operacional identificado no corpus, reforçando a plausibilidade dos padrões observados.

\begin{table}[H]
\centering
\caption{Critérios de priorização identificados no núcleo final}
\label{tab:criteriospriorizacao}
\small
\begin{tabular}{l r r}
\toprule
Critério & Registros & \% dos 104 \\
\midrule
Desempenho operacional & 93 & 89,4 \\
Informação e dados & 76 & 73,1 \\
Custo & 59 & 56,7 \\
Vida útil & 40 & 38,5 \\
Energia & 36 & 34,6 \\
Risco & 31 & 29,8 \\
Condição física & 27 & 26,0 \\
Manutenibilidade & 19 & 18,3 \\
Conforto & 17 & 16,3 \\
Segurança & 17 & 16,3 \\
Emissões de carbono & 15 & 14,4 \\
Resíduos & 13 & 12,5 \\
Satisfação do usuário & 9 & 8,7 \\
Qualidade de serviço & 6 & 5,8 \\
Água & 5 & 4,8 \\
\bottomrule
\end{tabular}

\vspace{2pt}
\begin{minipage}{0.75\linewidth}
\footnotesize Fonte: elaborado pelo autor a partir de \texttt{fontes/tabela26\_criterios\_nucleo\_final\_104.csv}.
\end{minipage}
\end{table}

A associação entre critérios e as cinco dimensões canônicas de sustentabilidade adotadas neste artigo -- ambiental, econômica, social, técnica-operacional e institucional -- está documentada, critério a critério, no dicionário de mapeamento produzido para esta revisão (\texttt{fontes/dicionario\_criterio\_dimensao\_etapa17.csv}), com justificativa individual para os dois casos de classificação menos evidente (risco e qualidade de serviço). Essa associação é insumo direto da matriz analítica apresentada na Seção~\ref{sec:matriz}.
